\chapter{The Instrument}
% Covers: outline B1-B3: the self-counting machine; the separability question; drift-immunity
% Chapter language (per Ch1 protocol): grammar. Arrow: lambda-calculus -> BNF.

\section{The machine that counts its own steps}
% Node-map: B1 (elaboration cost as THE observable; "heartbeat" = the machine's own tick,
% plain language, unit belongs to the apparatus; self-application licenses a COUNT --
% repeatable under stated procedure -- and no more; Ch8's wrap protected). EXP-2 previewed
% as apparatus practice (never trust a replay). APPARATUS-CITATION row placed HERE
% (bibliography-only; machine first named in this section). Gate: turn 514 (Ch3 S1, ~950 w).

The ground gave the argument its floor: atoms, addition, one identification, and names
priced as receipts. This chapter stands the instrument on that floor, and it speaks
grammar, as the map announced.

The instrument is a compiler: a machine that reads a text and rewrites it by rules
until no rule applies. Ours has one property that matters above all others, and the
argument leans on it everywhere: it terminates. On every lawful text, the rewriting
finishes. There is no text in this book's language on which the machine runs forever,
and so "the finished form of this text" is always a fact --- something the machine can
be asked for and will deliver.

One thing about this machine should be said plainly before the argument leans on it,
because a book that describes a device in prose can leave the impression the device
is a thought experiment. It is not. The instrument is a real program; it has been
built, it runs, and every reading this book reports was produced by running it and
checked by machine, not asserted on paper. Its source is public, and any reader who
doubts a number here can obtain the program and reproduce the number. What follows
describes the machine in words for the sake of the argument; the machine itself is
an artifact, not a thought experiment, and its existence is the one thing about it
this book does not ask the reader to take on trust. What its numbers \emph{mean} is
fenced with great care in the chapters ahead; \emph{that they were really computed}
is not fenced at all.

In grammar terms, a run is itself a text. One rule applied is one step:

\begin{center}
$T \;\longrightarrow\; T'$
\end{center}

and a run is the transcript of steps from the starting text to the form on which no
rule fires:

\begin{center}
$T \;\longrightarrow\; T_1 \;\longrightarrow\; T_2 \;\longrightarrow\; \cdots \;\longrightarrow\; T_n$
\end{center}

The machine ticks once per rule applied, and this book calls the machine's tick a
heartbeat, because the word is honest about whose pulse it is. A heartbeat is not a
second, not an energy, not a cost in any currency but the machine's own: it is one
rewrite, the atomic act of the instrument. The count of heartbeats in a run is the raw
observable of this entire book --- the needle's actual motion. And the reader should
notice that this count is nothing new in kind. A transcript is a text; its length is
read off it exactly as the roster lengths of the last chapter were read off theirs. The
instrument's observable is a parse report, and the ground's disciplines apply to it
unchanged: the count names the text it was read from, and the unit belongs to the
apparatus. When later chapters quote a count, it is quoted in heartbeats, read off the
run's transcript after the run --- the instrument's log, not the author's estimate.

Now the move that turns a compiler into an instrument. Write a description --- in the
machine's own language --- of the machine's act of compiling, and hand that text to the
machine. It compiles it. Nothing exotic occurs: a description of compiling is one more
lawful text, the machine terminates on lawful texts, and so it finishes and the
transcript has a length. The machine has counted its own steps --- more precisely, it
has counted the steps of compiling a description of its stepping --- and that number is
the self-measurement on which everything ahead is built.

Let the licence be stated exactly, because this is an edge on which cranks and mystics
both like to dance. What self-application licenses is a count: a number, produced by a
stated procedure, repeatable at will --- same text, same rules, same count, every run.
What it does not license is any form of self-knowledge. The machine does not become
aware of itself; it does not inspect its own soul; it reports a number, not an insight.
The recursion, moreover, has a floor: compiling a description of compiling is just
compiling. The tower does not run away, because each story of it is one more finite
text on which the machine terminates. Whether that tower of descriptions, pursued
honestly, closes into something --- whether the naming of namings comes back around ---
is a real question, and it is the closing chapter's question. This section spends none
of it. Here there is only the number, and the number is enough.

One piece of apparatus practice belongs beside the observable from the day it is
defined. The instrument, like any real instrument, keeps memories: finished forms of
texts it has compiled before, held so that ordinary work need not repeat itself. For
measurement, those memories are poison. A remembered count is a memory, not a
measurement, and the procedure of this book is that no count is quoted from memory:
before any measured run, the instrument's memory is emptied, and the run is made cold,
from the text, through every rewrite, to the finished form. Never trust a replay. The
full story of what this discipline caught --- and of the one count in the argument's
history that moved when a memory was disturbed --- belongs later in this chapter, where
the instrument's immunities are proved. Here it is stated as procedure, the way a
laboratory states that the balance is zeroed before each weighing.

A last word on credit. The particular machine used in this work is real, and it is
credited once, in the bibliography, as apparatus --- the way a paper names its
oscilloscope and moves on. In the prose it has no name and needs none: it is the
machine that counts its own steps, and every property of it the argument uses ---
termination, the tick, the transcript, the log --- has now been stated in the open.

The instrument is on the bench, then, and its needle moves. The next section asks the
instrument's first dangerous question: when the machine displays what it has counted,
does the displaying change the count?

\section{The separability question}
% Node-map: B2 whole (the K1-merged node WITH the two-instruments clause): ONE seam --
% does displaying the count change the count? Separable regime = proved harmless PER
% SITE by the static audit (EXP-5), re-runnable by any reader; entangled regime =
% EVIDENCED by EXP-1 (the heartbeat incident, told at last) -- evidenced NOT proved,
% not provable from inside, which IS the point. Repair (strong form) teased for S3.
% Discipline displayed as a fence. Gate: turn 515 (Ch3 S2, ~950 words).

An instrument that only measured would be safe, and an instrument that only displayed
would be harmless. This instrument does both with one body, and that is the danger. The
machine's reports must themselves be shown --- and a shown value is a text, and texts
are what the machine compiles, and compiling is what the measurement counts. In an
ordinary laboratory the needle and the dial are separate assemblies bolted to separate
frames. Here they share every part. So the question must be asked, and it deserves its
name: the separability question. When the machine displays what it has counted, does
the displaying change the count?

The suspicion is not paranoia; it is arithmetic. A displayed value occupies letters,
letters get rewritten, rewrites tick. On its face, every display is a finger resting on
the scale. The whole question is whether the finger presses.

It turns out there is not one answer. There are two regimes, and they are established
by two different kinds of instrument --- one a proof, one a scar.

The first regime is the separable one, and its instrument is a static audit. Many of
the machine's displayed values are read out and never read back: the value is shown,
and no rule of the running grammar ever takes the shown value as an input. That is a
property a reader can check by inspection, without running anything --- enumerate every
occurrence of the value's name across the texts; verify that each one is either the
value's own definition or an act of display; if so, the value feeds nothing, and
moving or removing its display cannot alter any count. For that value, at that site,
displaying is proved harmless. The proof is honest work but not deep work: it is an
afternoon with the texts and a pencil, and any reader may re-run it on any site at any
time. What must be respected is its scope. The audit proves one site at a time. It
licenses no generalization --- "displays are harmless" is not a theorem of this book
and never will be --- and every displayed value earns its own verdict or keeps the
status of a suspect.

Why such austerity, when site after site passes? Because of the second regime, which
is not a theorem but an event, and the argument carries its scar on file with a date.
The instrument once kept, inside a measured text, a line whose only job was display ---
a value shown for the reader's convenience, consumed, every inspection said, by
nothing. In a routine tidying, that line was moved out of the measured text into a
housekeeping text, the sort of relocation any engineer performs without a second
thought. And a count moved. One count, by one tick: a denominator that had read 2132
read 2131. Nothing had been miscounted, and that is precisely the horror of it --- the
machine counted correctly both times. The texts were different, because the displayed
line had been part of the text being measured. The machine counts its own elaboration;
the display lived inside what was elaborated; the dial, in this laboratory, has weight,
and the scale was weighing the dial.

Let the epistemic status of this be stated with the same exactness the licence got in
the last section. The entangled regime is evidenced, not proved. It is known the way
laboratories know things: it happened, on a stated day, to a stated count; it is on
file; the conditions reproduce it. And it could not have been proved from inside,
because the audit's whole method is to inspect the text --- and the incident consisted
of changing which text there was. The static instrument is blind in exactly the
dimension the event moved along. That is not a defect of the audit to be engineered
away; it is the boundary of what any inside inspection can certify, and it is why the
seam has two instruments and will keep both: the audit where the audit sees, the
incident's lesson where it cannot.

The discipline that survives, the working rule of this instrument for everything
ahead:

\begin{fence}
Separability is proved per site, never assumed. A displayed value whose harmlessness
has not been proved for its own site is treated as part of the measurement, and moving
it is treated as changing the experiment.
\end{fence}

One reassurance is owed before this section closes, and one debt. The count that moved
that day was not one the flagship reading rests on --- the bracket of Chapter 1
survived the incident untouched. But the argument does not ask the reader to find
comfort in that, because as it stood on the day, it was luck: the moved count might
have been load-bearing, and no audit would have said so in advance. Luck is not a
defense. The next section retires it --- it shows that the counts the flagship actually
rests on are immune to this entire class of disturbance, not by fortune but by proof,
and the proof has a shape worth the wait: everything that could drift, cancels.

\section{The immunities}
% Node-map: B3 (the strong form: three real-input forall-cancellation theorems -- the
% coupling's 18, the target's 5, the radius's 18 -- CARRIED, prose may say PROVED; the
% one-cancellation shape as algebra, one display; the surviving 18's two-by-nine root
% "checked, and checkable," NEVER "proved" -- PARTIAL kept honest). EXP-2 full story
% INCLUDING THE RETIREMENT (the fragile bare denominator withdrawn from every asserted
% claim). Closes: three immune inputs handed to Ch4 + one-clause fencepost-stop pointer
% (Ch4/5 owe the showing). Gate: turn 516 (Ch3 S3, ~1000 words, chapter closer).

The last section ended in dread, as it had to. One tick moved a count; the count that
moved happened not to matter; and happening is not a defense. If the readings ahead
rested on bare counts, any of them might be one relocation away from a different book.
This section is why they are not, and its answer is not care, and not luck. It is
proof.

The flagship reading rests on exactly three measured counts, and the reader has met
none of them yet by value. They can be named now, in the book's own terms: a coupling,
which the instrument measures as eighteen; a target, which it measures as five; and a
radius, which it measures as eighteen again. Chapter 4 will show where each comes from
and what each buys. What matters here is what they rest on --- because none of the
three is a bare count of heartbeats, and that is the entire architecture of their
safety.

Each of the three is a ratio, built so that the drift-prone bulk of the measurement
divides out. The shape is one line of algebra, and it deserves display:

\begin{center}
$\dfrac{18 \cdot M}{M} \;=\; 18 \qquad \text{for every } M > 0.$
\end{center}

Read it as the apparatus reads it. The coupling is measured as a numerator standing
over a product of the masses of everything that participated in the run --- and every
disturbance of the incident's class enters through that bulk. Heavier texts, longer
transcripts, a display relocated in or out, the machine elaborating more or less on the
way to its answer: all of it lands in $M$. And $M$ cancels. Not "cancels in the cases
we tried": for every $M$, as a theorem, proved about the actual quantity the reading
uses and not about a simplified stand-in. The same is proved, in the same real-input
form, for the target's five and for the radius's eighteen. Three counts, three
theorems, one shape: whatever the load, the reading is the reading. Everything that
could drift, cancels.

This is what retires the last section's luck, and it is worth saying precisely how.
The incident's lesson was that no inside inspection can promise a text will not
change. The immunity theorems do not promise that. They promise something better
placed: that when the text changes --- however it changes, by whatever tick --- the
quantities the flagship reads are unchanged, because change enters through the bulk
and the bulk divides out. The audit guards sites; the incident guards the audit's
pride; the cancellation guards the reading. Three instruments now, each honest about
its scope.

And the repair had a second half, which completes it. The count that moved that day
--- the bare denominator, the 2132 --- was not patched, or pinned, or apologized for.
It was retired. The argument withdrew that number from every claim it asserts: nowhere
ahead does this book state the denominator's value as a fact, because a bare count of
a machine's elaboration is exactly the kind of fact the incident proved fragile. What
the book asserts is the ratio, whose value survives the denominator's every mood. The
repair, in full: prove immunity for what the flagship reads, and withdraw assertion
from what it does not. After both halves, nothing load-bearing in this argument rests
on an unproved bare count --- the sentence Chapter 1 promised, earned here.

The whole affair, incident and repair alike, was conducted under the cold-run
discipline the first section stated as procedure, and this is its full story. The
incident was found because a measured run was made cold and its count compared against
the remembered one; a warm instrument would have replayed its memory and shown
nothing. The repair was verified the same way: memory emptied, every text recompiled
from the ground up, and the flagship reading demanded byte-for-byte identical before
and after. It was. Never trust a replay is not advice; it is how both halves of this
section's news were established.

One honest asterisk remains, and the book places it itself rather than leaving it for
a hostile reader. The cancellation leaves a survivor --- the numerator, the eighteen
--- and one may ask of it the question this chapter has taught: is the survivor a bare
count? Its form has been examined: the eighteen is twice nine, a structure with a
reason, not a happenstance total --- and the examination is repeatable by anyone. But
the book's own Chapter 2 distinguished an audited fact from a theorem, and the
distinction applies to the book's favorite number: the two-by-nine root of the
eighteen is checked, and checkable, and is not here called proved. The reader may hold
that asterisk against the argument; the argument prefers the asterisk to the inflation.

The instrument chapter is complete. On the bench: a machine that terminates and ticks;
a transcript whose length is the observable; a fence that prices every display; and
three measured inputs --- eighteen, five, eighteen --- immune by proof to the only
disturbance class ever seen to move a count. The counts are safe to spend. The next
chapter spends them: it is the reduction, where a count becomes a number, and where
the reader will watch the deepening procedure run its count to the fencepost of
Chapter 2 --- and stop there, as promised.
